\subsection{Why the henselized candidate does not settle L2}

The first run proposed the functor
$P(B)=\alpha_p(B^h_{(t)})$ on $\F_p[t]$-algebras, where the
superscript denotes henselization of the pair $(B,tB)$.
The question of its fpqc descent over arbitrary rings is separate
from the following obstruction. Even if that descent question has
a positive answer, this functor cannot refute L2.

\begin{proposition}\label{A:hensel-adic-failure}
\status{Proved}
For every prime $p$, the functor $P$ fails the strengthened
condition 4 on noetherian adic rings. More precisely, put
\[
 A_0=\F_p[t],\qquad R=A_0[[s]],\qquad
 A=R[\epsilon]/(\epsilon^2),\qquad I=(s).
\]
Then $A$ is noetherian and $I$-adically complete, but
\[
 P(A)\longrightarrow\varprojlim_{N\geq1}P(A/I^N)
\]
is not surjective.
\end{proposition}
\begin{proof}
Choose pairwise distinct monic irreducible polynomials
$q_0,q_1,\ldots\in\F_p[t]$, none equal to $t$, and set
$r_n=t/q_n(t)$. Write $D=(A_0)^h_{(t)}$ and $K=\F_p(t)$.
The ring $D$ is the henselian discrete valuation ring associated
with $(A_0)_{(t)}$, and its fraction field $K^h$ is a separable
algebraic extension of $K$. In $D$ there is a unique element $f_n$
such that
\[
 f_n^p-f_n=r_n,\qquad f_n\equiv0\pmod t.
\]
Here $q_n$ is a unit and the derivative of $X^p-X-r_n$ is $-1$,
so the henselian lifting criterion applies.

The classes of the $r_n$ are linearly independent in
$K/\{b^p-b:b\in K\}$ over $\F_p$. Indeed, any nonzero finite
linear combination has a pole of order one at one of the
irreducibles $q_n$. A rational function $b^p-b$ has at each pole
an order divisible by $p$, since the pole order of $b^p$ is
strictly greater than that of $b$.
It follows that
\[
 [K(f_0,\ldots,f_{N-1}):K]=p^N.                 \tag{A.2.1}
\]
For completeness, this compositum is a Galois extension whose
Galois group embeds in $\F_p^N$ by its translations of the roots.
If the image were proper, a nonzero linear functional vanishing
on it would give $\sum a_nf_n\in K$, and hence
$\sum a_nr_n=b^p-b$, a contradiction. Thus its group has order
$p^N$.

Henselization of pairs commutes with integral base change, in
particular with finite extensions and quotients
\cite[Lemma 15.12.7]{StacksHenselPairs}. Consequently
\[
 (A/s^N A)^h_{(t)}=D[s,\epsilon]/(s^N,\epsilon^2),
 \qquad
 A^h_{(t)}=R^h_{(t)}[\epsilon]/(\epsilon^2).
\]
The sections
\[
 z_N=\epsilon\sum_{0\leq n<N}f_ns^n
 \ \in\alpha_p\bigl((A/s^N A)^h_{(t)}\bigr)
\]
are compatible. Their $p$th powers vanish for every prime $p$.
Suppose that they were induced by $z\in P(A)$.
The ring $R^h_{(t)}$ is reduced: it is a filtered colimit of
\'etale algebras over the reduced ring $R$, and these algebras
are reduced. This is the reducedness argument in
\cite[Lemma 15.46.4]{StacksHenselReduced}, which applies verbatim
to a pair. Thus $z=\epsilon a$ for some $a\in R^h_{(t)}$.

There is a canonical homomorphism
\[
 R^h_{(t)}\longrightarrow
 \varprojlim_N (R^h_{(t)}/s^N R^h_{(t)})=D[[s]].
\]
The image of $a$ must be $\sum_{n\geq0}f_ns^n$.
By the construction of henselization, $a$ comes from one
finitely presented \'etale $R$-algebra $C$ equipped with
$C\to R^h_{(t)}$; see
\cite[Lemma 15.12.1]{StacksHenselPairs}.
The induced map $C/sC\to D\to K^h$ has, after extension from
$A_0$ to $K$, image a finite separable field $L/K$: a finitely
presented \'etale algebra over a field is a finite product of
finite separable extensions, and a map to a field selects one
factor.

Formal \'etaleness now gives a unique compatible system of
$R$-algebra maps $C\to L[s]/(s^N)$ lifting $C/sC\to L$,
and therefore a map $C\to L[[s]]$.
After extension to $K^h[[s]]$, this map and the original map
$C\to R^h_{(t)}\to D[[s]]$ have the same reduction modulo $s$.
Formal \'etale uniqueness, applied successively modulo $s^N$,
shows that they agree. Every coefficient of the image of $a$
therefore lies in the one finite extension $L/K$. This contradicts
(A.2.1), since all the $f_n$ would lie in $L$.
The displayed compatible system has no preimage, as asserted.
\end{proof}

The argument uses formal \'etaleness of the single finite
presentation $C$. It does not posit a map from the whole
henselization $R^h_{(t)}$ to $K[[s]]$: such a map is not supplied
by the universal property of henselian pairs when $t$ becomes a
unit. This distinction is essential to the coefficient argument.
The proposition eliminates the proposed candidate under the
strengthened axiom; it does not prove or refute L2 itself.
